% !TeX root = ../main.tex
\section{Full-tracker baseline}
\label{sec:tracker-overview}
The tracker combines three silicon technologies in a common cylindrical layout:
pixels close to the beam pipe, short-strip strixels at intermediate radii, and
long-strip stereo pairs at the outer radii. This edition fixes the selected
layouts, local supports and service reservations as a baseline for technical
review. The dimensions below are nODD design choices. They describe the selected
models, including their engineering approximations; they do not constitute
construction approval or a demonstrated tracking-performance specification.

The pixel system supplies two-coordinate measurements for vertexing and forward
tracking. The short-strip sensor also measures two coordinates, using a fine
transverse pitch and a coarser second coordinate. The outer system uses two
separate one-dimensional strip sensors at a small relative stereo angle. Its
larger lever arm is intended to constrain the track curvature while controlling
channel count. The long-strip pair needs a trajectory model to combine its
separated measurements; two sensor faces are not interchangeable with a single
two-coordinate plane.

\subsection{Layout and coordinate conventions}
The beam direction defines $z$, and $r=\sqrt{x^2+y^2}$ is the transverse radius.
The nominal interaction point is at the origin. The selected architecture has
ten barrel assemblies: four pixel, four short-strip and two long-strip layers.
Each signed end contains ten pixel discs, six short-strip discs and six
long-strip discs, giving 22 disc assemblies per end. A long-strip assembly contains both faces of every stereo
pair. \Cref{fig:tracker-layout,tab:tracker-inventory} distinguish assembly counts
from sensitive-element counts.

\begin{figure}[tbp]
\centering\includegraphics[width=\linewidth]{tracker-layout}
\caption{Selected tracker layer and disc datums in the $r$--$z$ plane. The
complete barrel and positive discs are shown; the negative end has corresponding
signed placements. The lines mark nominal layers and active annuli, rather than
all staggered sensor planes or passive volumes. Dashed rays are straight
trajectories from the nominal interaction point, not coverage predictions in
the field. The red origin interval illustrates the luminous-region hypothesis.}
\label{fig:tracker-layout}
\end{figure}

\begin{table}[tbp]
\centering\small
\caption{Selected tracker inventory. Endcap totals include both ends. Pixel
chips are active readout patches; strixels are complete two-coordinate sensors;
long-strip faces are separate one-dimensional sensors. The maintained full
tracker contains \num{46250} sensitive elements.}
\label{tab:tracker-inventory}
\begin{tabular}{@{}lrrp{58mm}@{}}\toprule
Assembly & Layers/discs & Modules/pairs & Sensitive inventory\\\midrule
Pixel barrel & 4 & \num{3050} & \num{6794} chip patches\\
Pixel endcaps & $2\times10$ & \num{2616} & \num{5616} chip patches\\
Short-strip barrel & 4 & \num{8064} & \num{8064} strixels\\
Short-strip endcaps & $2\times6$ & \num{4320} & \num{4320} strixels\\
Long-strip barrel & 2 & \num{6984} pairs & \num{13968} strip faces\\
Long-strip endcaps & $2\times6$ & \num{3744} pairs & \num{7488} strip faces\\\bottomrule
\end{tabular}
\end{table}

Pixel barrel datums are at $r=31.5,57,96,182\,\unit{\milli\metre}$, with nominal
active extent $|z|\leq\qty{550}{\milli\metre}$. Short-strip stave anchors lie at
$r=260,340,480,660\,\unit{\milli\metre}$ with nominal active extent
$|z|\leq\qty{1200}{\milli\metre}$. Long-strip anchors lie at
$r=840,1060\,\unit{\milli\metre}$, with an inherited nominal active half-length
of \qty{1287.33}{\milli\metre}. These radii are datums, not average sensor radii:
tilt, module elevation, passive stacks and stereo separation give each sensitive
element its own position and normal.

Disc support datums extend to \qty{3070}{\milli\metre} for pixels,
\qty{3080}{\milli\metre} for short strips and \qty{3120}{\milli\metre} for long
strips. The first short-strip disc is at \qty{1335}{\milli\metre} and the first
long-strip disc at \qty{1470}{\milli\metre}; both were moved outward to clear
their barrel collectors. The associated old-reference coverage losses are
retained in the subsystem sections. The geometry is therefore a coherent
review layout, with integration and transition-performance questions still open.

The short-strip barrel uses a common $+\qty{15}{\degree}$ azimuthal tilt and the
long-strip barrel a common $+\qty{12}{\degree}$ tilt. The complete cold structure
rotates with its sensors, tubes and local services. This removes the
alternating-radius stave lanes of the earlier short-strip layout and makes the
local geometry covariant under a one-stave rotation. Different layer
populations and sector services still produce an azimuth-dependent material
distribution. Neither this symmetry nor geometric overlap establishes uniform
resolution. In the endcaps, module axes follow the local polar frame rather
than global $x$--$y$ directions. Proper rotations define the negative-end strip
frames and their identifiers.

\subsection{Coverage hypotheses and measurement model}
The architecture targets charged-particle tracking to $|\eta|<4$. The pixel
active apertures retain the selected ring populations, whose original design
screen considered first outward trajectories with
$\pt\geq\qty{1}{\giga\electronvolt}$, a luminous interval
$z_0\in[-150,150]\,\unit{\milli\metre}$ and the stated uniform axial-field
hypotheses up to \qty{4}{\tesla}. This interval is a design-screen envelope,
not a measured beam distribution. It permits progressively larger pixel active
apertures at the outer stations while retaining their common mechanical plate.
It does not imply that every strip subsystem independently covers $|\eta|<4$.

The retained strip screens sample straight tracks and vacuum helices at selected
momenta, fields, vertices and azimuths. They use finite sensor planes and retain
the original ideal-cylinder or annulus denominators. They are neither proofs
for every intermediate momentum nor simulations of scattering, energy loss,
dead channels or track reconstruction. Nonzero misses remain at barrel ends and
shifted first discs. A finite Fatras truth-tracking comparison uses origin muons at
$\pt=1,10,100\,\unit{\giga\electronvolt}$, $|\eta|\leq4$ and a constant
\qty{3}{\tesla} axial field, with geometric digitization as the primary
response. About \qty{98.2}{\percent} of generated tracks have at least ten
physical measurements. This integrated fraction retains forward gaps; it is
not a minimum-hit guarantee throughout the acceptance. The selected geometry
also has a low-momentum resolution tradeoff. Detailed residuals, pulls and
material/response assumptions belong to the accompanying validation evidence.

Pixel readout uses square \qty{50}{\micro\metre} cells. The short-strip fixture
uses \qty{75}{\micro\metre} by \qty{0.5}{\milli\metre} cells, with no added
stereo pair. Long-strip sensors use \qty{80}{\micro\metre} one-dimensional
pitch and physical axes at $\pm\qty{20}{\milli\radian}$ about the nominal
module direction. These are segmentation and geometry choices. Charge sharing,
Lorentz drift, alignment covariance and calibrated spatial resolution remain
separate response-model requirements.

\subsection{Support hierarchy and shared services}
The load path is sensor and electronics, local thermal contact, structural
sandwich, bearing or petal contacts, then subsystem frames and the global tracker
support. In the long-strip system the cooled sandwich is itself the stave or
petal; a second hidden carrier is not added between the stereo sensors. The
pixel endcaps share a closed carrier with longitudinal rails. Strip barrels use
bearing rings, and strip discs use removable petals located by constrained and
sliding contacts. These local constructions define load transfer and thermal
motion concepts. A qualified global cage, installation tooling and connections
to the surrounding detector are still required.

All three systems route local LV, HV and signal services to accessible end
boards or rim interfaces, and then to longitudinal trunks. Embedded titanium
tubes carry the provisional evaporative CO$_2$ circuits. Supply, return, cable
strain relief and sensor mounting need independent constraints so service loads
do not act through silicon bonds. Local loops and common manifolds must also
permit pressure testing and subsystem assembly before installation.

\begin{figure}[tbp]
\centering\includegraphics[width=\linewidth]{tracker-services}
\caption{Positive-side service reservations at the selected fixed radii. Filled
regions show barrel collection bays, disc fans and longitudinal trunk envelopes;
they are not filled solids or a cable-installation model. The short-strip
envelope is $r=710$--$783$ mm and the long-strip envelope
$r=1144$--$1169$ mm. Pixel services occupy their separate inner corridor.
Individual fittings, bend sweeps and combined subsystem clearance require review.}
\label{fig:tracker-services}
\end{figure}

The inner pixel trunk reserves $r=190$--$234\,\unit{\milli\metre}$ before the
rear handoff. The short-strip trunk reserves $r=710$--$783\,\unit{\milli\metre}$
from $|z|=1310$ to \qty{3500}{\milli\metre}; the long-strip trunk reserves
$r=1144$--$1169\,\unit{\milli\metre}$ from $|z|=1435$ to
\qty{3500}{\milli\metre}. These fixed envelopes are not expanded when a capacity
screen fails. Their separated radial bands alone do not prove installation
clearance at collection bays, frames or the rear transitions.

Detailed models contain explicit local sensors, foams, skins, pickups, pipes and
selected mounting components. Remote services additionally use effective cells
that conserve constituent volumes and residual void. Gross cable area remains
a separate packing budget. A combined model must replace the barrel-only strip
service cells with the corresponding cumulative barrel-plus-endcap inventory;
overlaying independent compacts would allocate the same corridor and material
twice. For the same reason, their standalone model masses must not be summed
as a tracker mass or material budget. The maintained integrated exporter reconciles the strip service inventories.
Global pixel service capacity and fresh integrated directional material maps
remain separate review deliverables.
